\documentclass[main.tex]{subfiles}


\begin{document}

%from pagenumber to up
% \phantomsection 
% \hypertarget{toc}{}

\pagestyle{empty}

\begin{NoHyper} % отключить клик-ссылки

% номер секции вручную
\setcounter{section}{33
}
\addtocounter{section}{-1} 

\section{Магнитное поле кольца с током}


Пусть имеется провод в форме кольца, через который протекает электрический ток. На рис.\,\ref{fig:p149} показана картина линий магнитного поля кольца с током.

\makeatletter

\pgfkeys{/pgf/decoration/.cd,
         start color/.store in =\startcolor,
         end color/.store in   =\endcolor
}

\pgfdeclaredecoration{width and color change}{initial}{
 \state{initial}[width=0pt, next state=line, persistent precomputation={%
   \pgfmathdivide{50}{\pgfdecoratedpathlength}%
   \let\increment=\pgfmathresult%
   \def\x{0}%
 }]{}
 \state{line}[width=.5pt,   persistent postcomputation={%
     \pgfmathadd@{\x}{\increment}%
     \let\x=\pgfmathresult%
   }]{%
   \pgfsetlinewidth{\x/150*0.035pt+\pgflinewidth}%
   \pgfsetarrows{-}%
   \pgfpathmoveto{\pgfpointorigin}%
   \pgfpathlineto{\pgfqpoint{.75pt}{0pt}}%
   \pgfsetstrokecolor{\endcolor!\x!\startcolor}%
   \pgfusepath{stroke}%
 }
 \state{final}{%
   \pgfsetlinewidth{\pgflinewidth}%
   \pgfpathmoveto{\pgfpointorigin}%
   \color{\endcolor!\x!\startcolor}%
   \pgfusepath{stroke}% 
 }
}

\makeatother


    \begin{figure}[H]
    \centering

    \begin{tikzpicture}[font=\small,            line cap=round,                             >={Stealth[bend,inset=0pt,round,length=3mm, angle'=20]},vec/.style={>={Stealth[inset=0pt,round,length=3.5mm, angle'=20]}}]


\def\lm{1} 
\def\xr{.25}


%%%%%%%%%%%%%%%%%%%%%%%%%%%%%
%%%%%%%%%%%%%%%%%%%%%%%%%%%%%

    
%lines left
%1
\def\ds{.4}
\path[yshift=.75cm,->] (0.25,0) 
arc [start angle=-90, end angle=0, x radius={\ds*1cm}, y radius={\ds*1cm-.02cm}]
arc [start angle=0, end angle=90, x radius={\ds*1cm}, y radius={\ds*1cm+.05cm}]
coordinate (e1)
;

\draw[gray,yshift=.75cm] (e1) 
arc [start angle=90, end angle=180, x radius={\ds*1cm}, y radius={\ds*1cm+.05cm}] 
arc [start angle=180, end angle=270, x radius={\ds*1cm}, y radius={\ds*1cm-.02cm}] 
;


%2
\def\df{.9}
\path[yshift=.5cm,->] (0.25,0) 
arc [start angle=-90, end angle=0, x radius={\df*1cm}, y radius={\df*1cm-.1cm}]
arc [start angle=0, end angle=90, x radius={\df*1cm}, y radius={\df*1cm+.1cm}]
coordinate (e1)
;

\draw[gray,yshift=.5cm] (e1) 
arc [start angle=90, end angle=180, x radius={\df*1cm}, y radius={\df*1cm+.05cm}] 
arc [start angle=180, end angle=270, x radius={\df*1cm}, y radius={\df*1cm-.05cm}] 
;


%3
\def\dg{1.5}
\path[yshift=.25cm,->] (0.25,0) 
arc [start angle=-90, end angle=0, x radius={\dg*1cm}, y radius={\dg*1cm-.1cm}]
arc [start angle=0, end angle=90, x radius={\dg*1cm}, y radius={\dg*1cm+.1cm}]
coordinate (e1)
;

\draw[gray,yshift=.5cm] (e1) 
arc [start angle=90, end angle=180, x radius={\dg*1cm}, y radius={\dg*1cm+.05cm}] 
arc [start angle=180, end angle=270, x radius={\dg*1cm}, y radius={\dg*1cm-.05cm}] 
;

%reverse
\begin{scope}[yscale=-1,gray]

%1
\path[yshift=.75cm,->] (0.25,0) 
arc [start angle=-90, end angle=0, x radius={\ds*1cm}, y radius={\ds*1cm-.02cm}]
arc [start angle=0, end angle=90, x radius={\ds*1cm}, y radius={\ds*1cm+.05cm}]
coordinate (e1)
;

\draw[yshift=.75cm] (e1) 
arc [start angle=90, end angle=180, x radius={\ds*1cm}, y radius={\ds*1cm+.05cm}] 
arc [start angle=180, end angle=270, x radius={\ds*1cm}, y radius={\ds*1cm-.02cm}] 
;


%2
\def\df{.9}
\path[yshift=.5cm,->] (0.25,0) 
arc [start angle=-90, end angle=0, x radius={\df*1cm}, y radius={\df*1cm-.1cm}]
arc [start angle=0, end angle=90, x radius={\df*1cm}, y radius={\df*1cm+.1cm}]
coordinate (e1)
;

\draw[yshift=.5cm] (e1) 
arc [start angle=90, end angle=180, x radius={\df*1cm}, y radius={\df*1cm+.05cm}] 
arc [start angle=180, end angle=270, x radius={\df*1cm}, y radius={\df*1cm-.05cm}] 
;


%3
\def\dg{1.5}
\path[yshift=.25cm,->] (0.25,0) 
arc [start angle=-90, end angle=0, x radius={\dg*1cm}, y radius={\dg*1cm-.1cm}]
arc [start angle=0, end angle=90, x radius={\dg*1cm}, y radius={\dg*1cm+.1cm}]
coordinate (e1)
;

\draw[yshift=.5cm] (e1) 
arc [start angle=90, end angle=180, x radius={\dg*1cm}, y radius={\dg*1cm+.05cm}] 
arc [start angle=180, end angle=270, x radius={\dg*1cm}, y radius={\dg*1cm-.05cm}] 
;

\draw[decoration={markings,mark=at position {1/2+1.5/20} with {\arrow{<}}},postaction={decorate}] (.25,0) --++ (-2,0);

\end{scope}


%%%%%%%%%%%%%%%%%%%%%%%%%%%%%
%ring
\begin{scope}
    
\draw[line width=.4pt, decoration={width and color change,start color=brown, end color=brown}, decorate] (0.5,{.0pt})
arc [start angle={0}, end angle={180}, x radius={\xr*1cm}, y radius=1cm] 
% arc [start angle=270, end angle=180, x radius=.5cm, y radius=1cm] 
;

\draw[line width=.4pt, decoration={width and color change,start color=brown, end color=brown}, decorate] (0.5,{-.0pt})
% arc [start angle={90}, end angle={180}, x radius=.5cm, y radius=1cm] 
arc [start angle=360, end angle=180, x radius={\xr*1cm}, y radius=1cm] 
;

%arrows
\path[red,->,thin] (-0.2,0)
% arc [start angle={90}, end angle={180}, x radius=.5cm, y radius=1cm] 
arc [start angle=180, end angle=190, x radius={\xr*1cm}, y radius=1cm] coordinate (a) 
;

\draw[red,->,thin] (a)
% arc [start angle={90}, end angle={180}, x radius=.5cm, y radius=1cm] 
arc [start angle=190, end angle=260, x radius={\xr*1cm}, y radius=1cm] node[black,below right,shift={(-.16,.07)}] {$I$}
;

\end{scope}


%%%%%%%%%%%%%%%%%%%%%%%%%%%%%
%%%%%%%%%%%%%%%%%%%%%%%%%%%%%


%lines right
%1
\draw[gray,yshift=.75cm,->] (0.25,0) 
arc [start angle=-90, end angle=0, x radius={\ds*1cm}, y radius={\ds*1cm-.02cm}]
arc [start angle=0, end angle=90, x radius={\ds*1cm}, y radius={\ds*1cm+.05cm}]
coordinate (e1)
;


%2
\draw[gray,yshift=.5cm,->] (0.25,0) 
arc [start angle=-90, end angle=0, x radius={\df*1cm}, y radius={\df*1cm-.1cm}]
arc [start angle=0, end angle=90, x radius={\df*1cm}, y radius={\df*1cm+.1cm}]
coordinate (e1)
;


%3
\def\dg{1.5}
\draw[gray,yshift=.25cm,->] (0.25,0) 
arc [start angle=-90, end angle=0, x radius={\dg*1cm}, y radius={\dg*1cm-.1cm}]
arc [start angle=0, end angle=90, x radius={\dg*1cm}, y radius={\dg*1cm+.1cm}]
coordinate (e1)
;


%reverse
\begin{scope}[yscale=-1,gray]

%1
\draw[yshift=.75cm,->] (0.25,0) 
arc [start angle=-90, end angle=0, x radius={\ds*1cm}, y radius={\ds*1cm-.02cm}]
arc [start angle=0, end angle=90, x radius={\ds*1cm}, y radius={\ds*1cm+.05cm}]
coordinate (e1)
;


%2
\draw[yshift=.5cm,->] (0.25,0) 
arc [start angle=-90, end angle=0, x radius={\df*1cm}, y radius={\df*1cm-.1cm}]
arc [start angle=0, end angle=90, x radius={\df*1cm}, y radius={\df*1cm+.1cm}]
coordinate (e1)
;


%3
\def\dg{1.5}
\draw[yshift=.25cm,->] (0.25,0) 
arc [start angle=-90, end angle=0, x radius={\dg*1cm}, y radius={\dg*1cm-.1cm}]
arc [start angle=0, end angle=90, x radius={\dg*1cm}, y radius={\dg*1cm+.1cm}]
coordinate (e1)
;

\draw[decoration={markings,mark=at position {1/2+1.5/20} with {\arrow{>}}},postaction={decorate}] (.25,0) --++ (2,0);

\end{scope}

    
    \end{tikzpicture}
\caption{Линии магнитного поля кольцевого провода с током}
\label{fig:p149}
\end{figure}


% Через любую точку пространства вне кольца можно провести соответствующую линию его магнитного поля. С удалением от кольца густота линий поля убывает (убывает индукция магнитного поля).

Линии поля кольца как бы выходят из него со одной стороны и входят в него с другой стороны (внутри кольца линии замыкаются).

Для определения того, с какой стороны из кольца выходят линии его магнитного поля, можно воспользоваться следующим правилом.
%
\begin{law}
    \textbf{Правило правой руки.} Направление линий магнитного поля \emph{внутри} кольца укажет большой палец правой руки, если остальные согнутые пальцы (охватывающие кольцо) расположены так, что они указывают направление тока кольца.
\end{law}
%
Рисунок~\ref{fig:p150} иллюстрирует применение этого правила.  

    \begin{figure}[H]
    \centering
\begin{asy}
settings.outformat="png";
// settings.prc=true;
//settings.render=10;

import three;
import texcolors;
import x11colors;
import solids;
size(10cm,0);
currentprojection =
orthographic((5,3,1.5));


//styles
///////////////////////////////////////////
DefaultHead3.size=new real(pen p=currentpen) {return 3mm;};
defaultpen(fontsize(11pt));
defaultpen(linewidth(0.4pt));
pen thick=black+2*linewidth(currentpen);
/////////////////////////////////////////


//draw(surface(path3D), lightgray+opacity(.5));
//draw(path3D, Chocolate+thick,L=Label("$S$",EndPoint,N,black));

//axes
//draw(O--2.5X,L=Label("$x$",EndPoint,WSW));
//draw(O--2.3Y,L=Label("$y$",EndPoint,ESE));
//draw(O--3Z,L=Label("$\vec n$",EndPoint,W),Arrow3(angle=10));
//draw(O--rotate(-30,X)*4Z,Green+thick,L=Label("$\vec B$",EndPoint,E,black),Arrow3(3.5mm,angle=10,emissive(Green)));

//draw("$\alpha$",reverse(arc(O,0.8*Z,rotate(-30,X)*0.8*Z)),N+0.3E);


//curr
//cylinder1
path3 c2=circle((-.8,2.2,0),0.1,Z);
revolution sur1=revolution((0,0,0),c2,Y,0,360);
draw(surface(sur1),orange);

path3 arc1=arc((0,2,0),1.1,70,0,135,0,CCW);
draw(arc1,red+thick,L=Label("$I$",MidPoint,W,white),Arrow3(3.5mm,angle=10,emissive(red)));


////////////////////////////////////////
//plates
real a=2;
path3 pl=((0,-a,-a)--(0,-a,a)--(0,a,a)--(0,a,-a)--cycle); 
draw(shift(-1,0,0)*surface(pl), lightgray+opacity(1));
draw(shift(-2,0,0)*surface(pl), lightgray+opacity(1));

//cyl
triple pO=(-1.5,2,-2);
revolution cyl=cylinder(pO,.5,4,Z);
draw(surface(cyl),lightgray);

triple pO=(-1.5,-2,-2);
revolution cyl=cylinder(pO,.5,4,Z);
draw(surface(cyl),lightgray);

triple pO=(-1.5,-2,2);
revolution cyl=cylinder(pO,.5,4,Y);
draw(surface(cyl),lightgray);

triple pO=(-1.5,-2,-2);
revolution cyl=cylinder(pO,.5,4,Y);
draw(surface(cyl),lightgray);

//sphere
draw(shift(-1.5,2,-2)*scale3(.5)*unitsphere,lightgray);
draw(shift(-1.5,2,2)*scale3(.5)*unitsphere,lightgray);
draw(shift(-1.5,-2,-2)*scale3(.5)*unitsphere,lightgray);
draw(shift(-1.5,-2,2)*scale3(.5)*unitsphere,lightgray);


//fingers
for (int n = 0; n <= 3; ++n) {

triple f1=(-1.5,a,a);
revolution cyl=cylinder(f1,.5,((5/8)*2a),Z);
draw(rotate(-60,(-1.5,a,a),(-1.5,-a,a))*shift(0,(-1-.333)*n,0)*surface(cyl),lightgray);
  
triple f2=(shift((5/8)*2a*Sin(60),0,(5/8)*2a*Cos(60))*f1);
revolution cyl=cylinder(f2,.5,((3/8)*2a),Z);
draw(rotate(-60-70,f2,shift(0,-1,0)*f2)*shift(0,(-1-.333)*n,0)*surface(cyl),lightgray);
draw(shift(f2)*shift(0,(-1-.333)*n,0)*scale3(.5)*unitsphere,lightgray);
  
//dot(shift((3/8)*2a*Sin(60+70),0,(3/8)*2a*Cos(60+70))*f2,red);
triple f3=(shift((3/8)*2a*Sin(60+70),0,(3/8)*2a*Cos(60+70))*f2);
revolution cyl=cylinder(f3,.5,((2/8)*2a),Z);
draw(rotate(-60-70-50,f3,shift(0,-1,0)*f3)*shift(0,(-1-.333)*n,0)*surface(cyl),lightgray);
draw(shift(f3)*shift(0,(-1-.333)*n,0)*scale3(.5)*unitsphere,lightgray);
  
triple f4=(shift((2/8)*2a*Sin(60+70+50),0,(2/8)*2a*Cos(60+70+50))*f3);
draw(shift(f4)*shift(0,(-1-.333)*n,0)*scale3(.5)*unitsphere,lightgray);

}


//big fing
triple pO=(-1.5,2,-2);
revolution cyl=cylinder(pO,.5,2,Y);
draw(rotate(-10,pO,shift(-1,0,0)*pO)*surface(cyl),lightgray);

triple bf=(shift(0,(4/8)*2a*Cos(10),(4/8)*2a*Sin(10))*pO);
revolution cyl=cylinder(bf,.5,((2.5/8)*2a),Y);
draw(surface(cyl),lightgray);
draw(shift(bf)*scale3(.5)*unitsphere,lightgray);
draw(shift(0,(2.5/8)*2a,0)*shift(bf)*scale3(.5)*unitsphere,lightgray);




\end{asy}
\caption{Правило правой руки}
\label{fig:p150}
\end{figure}


Соотнося рисунки~\ref{fig:p149} и~\ref{fig:p150}, можно видеть, что большой палец на рис.\,\ref{fig:p150} указывает направление линий магнитного поля внутри кольца.


















\end{NoHyper}
\end{document}
